% ============================================================
% Project: The Philosophy of Viability
% Book: Book I — Ontology of Viability
% Book ID: VS.B01
% File: section_06_plotinus.tex
% Section: Chapter 08, Section 06 — Plotinus
% Purpose: Historical reconstruction and critical comparison
% ============================================================

\section{Plotinus: The One, Intellect, Soul, and the Ordered Manifestation of Reality}
\label{sec:ch08-plotinus}

\subsection*{Scope and Comparative Orientation}

This section examines Plotinus as a major architect of late-antique
metaphysics. Its purpose is not to assimilate his philosophy to the
ontology developed in this book, nor to treat Neoplatonism as an
anticipation of the Philosophy of Viability. It is to reconstruct
Plotinus's principal ontological commitments, identify the explanatory
relations among them, and determine where his account provides a
genuine point of comparison, where it differs in kind, and where the
comparison must remain open.

The inquiry concentrates on five interconnected questions:

\begin{enumerate}
    \item What is the status of the One as the first principle?
    \item How are Intellect and Soul related to the One and to one
    another?
    \item How does Plotinus account for the ordered appearance of
    sensible reality?
    \item What roles do unity, multiplicity, form, matter, and
    participation play in this account?
    \item What can be learned from Plotinus without importing his
    metaphysical hierarchy into the ontology of viability?
\end{enumerate}

The principal textual basis is the \textit{Enneads}, especially
V.1, V.2, V.3, V.4, V.5, V.9, VI.7, VI.8, VI.9, III.8, IV.8,
and I.8. These treatises must be read together. Porphyry's thematic
arrangement is indispensable for locating the discussions, but it
should not be mistaken for a transparent reproduction of the order
in which Plotinus developed his thought.

\subsection{Historical and Textual Location}
\label{subsec:ch08-plotinus-historical}

\phantomsection\label{para:VS.B01.P01.C08.S06.SS01.P001}
Plotinus (204/5--270 CE) stands within the Platonic tradition and
develops a comprehensive metaphysical account whose central principles
are the One or Good, Intellect, and Soul. His philosophy is neither
merely a commentary on Plato nor a simple restatement of earlier
Platonism. It is a systematic reconstruction of Platonic materials
under the pressure of questions raised by Aristotle, the Stoics, and
other philosophical currents of late antiquity.

\phantomsection\label{para:VS.B01.P01.C08.S06.SS01.P002}
The surviving body of Plotinus's writings is transmitted through
Porphyry's edition of the \textit{Enneads}. Porphyry arranged the
treatises into six groups of nine. This arrangement is thematic rather
than chronological. It is therefore useful for comparative study but
does not, by itself, establish the developmental order of Plotinus's
arguments. The treatises frequently revisit the same problems from
different directions, and a responsible reconstruction must compare
passages across the collection rather than extract a complete doctrine
from a single treatise.

\phantomsection\label{para:VS.B01.P01.C08.S06.SS01.P003}
The historical relation to Aristotle is particularly important for
the present chapter. Plotinus engages Aristotelian questions about
substance, form, matter, actuality, intellect, and the kinds of being.
He does not simply reject Aristotle. He appropriates, transforms, and
criticizes Aristotelian analyses within a metaphysical architecture
whose first principle exceeds the order of being and intelligibility.
The comparison with Aristotle in Section~05 should therefore be
understood as a transition between distinct explanatory architectures,
not as a linear development from one system into the other.

\subsection{The Problem of the First Principle}
\label{subsec:ch08-plotinus-first-principle}

\phantomsection\label{para:VS.B01.P01.C08.S06.SS02.P001}
The central problem of Plotinus's metaphysics is how a reality that is
multiple, differentiated, intelligible, and sensible can depend upon
a first principle that is absolutely simple. The first principle cannot
be one being among other beings, since it would then belong to the
order whose existence it is supposed to explain. Nor can it be a
composite entity, since its own composition would require an account
of the principles that compose it. Plotinus therefore identifies the
first principle as the One, also called the Good.

\phantomsection\label{para:VS.B01.P01.C08.S06.SS02.P002}
The One is not a numerical unit within a plurality. Numerical unity
presupposes a domain in which units can be counted and distinguished.
The One, by contrast, is prior to the differentiation through which
countable beings become possible. It is not one item in a collection,
nor the largest member of a class, nor a totality assembled from
parts. Its unity is intended to be absolute simplicity rather than
numerical singularity.

\phantomsection\label{para:VS.B01.P01.C08.S06.SS02.P003}
This distinction is essential. If the One were treated as a singular
object, it would already be situated within a field of possible
objects and relations. Plotinus's claim is more radical: the One is
the principle on which the possibility of determinate being depends,
while it is not itself adequately characterized by the determinations
that belong to beings.

\phantomsection\label{para:VS.B01.P01.C08.S06.SS02.P004}
The One is consequently described through a language of transcendence
and negation. It is beyond being, beyond predication, and beyond
ordinary intellection. Such expressions do not mean that the One is
simply nonexistent. They indicate that the categories through which
determinate beings are understood cannot be transferred without
qualification to the first principle. Plotinus's language is therefore
not a license to turn the One into an indeterminate object. It marks
the limit of the conceptual forms that apply to determinate reality.

\phantomsection\label{para:VS.B01.P01.C08.S06.SS02.P005}
The identification of the One with the Good adds a second dimension.
The One is not merely the first explanatory term in a causal sequence.
It is also that toward which all things are oriented as toward their
good. This orientation is not necessarily a conscious intention in
every being. It is a metaphysical structure of dependence and
aspiration. The Good is not one good among others; it is the
transcendent principle by reference to which the goodness of other
things is understood.

\subsection{Intellect: Being, Form, and Intelligible Multiplicity}
\label{subsec:ch08-plotinus-intellect}

\phantomsection\label{para:VS.B01.P01.C08.S06.SS03.P001}
The second principle is Intellect (\textit{Nous}). Intellect is not
the One, because intellection involves a structure of intelligibility:
that which thinks and that which is thought are distinguishable, even
when they belong to one unified activity. The One, being absolutely
simple, cannot be an intellect in this sense. Intellect therefore
marks the emergence of determinate intelligibility and articulated
multiplicity.

\phantomsection\label{para:VS.B01.P01.C08.S06.SS03.P002}
In Intellect, being and intelligibility are inseparable. The Forms are
not external objects that Intellect subsequently observes. They
constitute the intelligible content of Intellect itself. Intellect is
therefore a living unity of Forms, not an empty faculty that receives
representations from an independent domain. Its multiplicity is
internal and intelligible: the Forms are distinct, yet they belong to
one articulated intelligible life.

\phantomsection\label{para:VS.B01.P01.C08.S06.SS03.P003}
This account differs from any model in which unity is achieved by
eliminating distinctions. Intelligible unity does not mean that the
Forms become indistinguishable. Nor does multiplicity necessarily
imply external separation. Plotinus seeks to explain how determinate
difference can be present within a unity that is not assembled from
independent parts.

\phantomsection\label{para:VS.B01.P01.C08.S06.SS03.P004}
The relation between the One and Intellect is often described through
the language of procession. Intellect proceeds from the One, but this
procession is not a temporal event in which the One first exists alone
and later produces another entity. It expresses an ontological
dependence. Intellect is dependent upon the One while possessing a
determinate mode of being that the One itself does not possess.

\phantomsection\label{para:VS.B01.P01.C08.S06.SS03.P005}
Intellect also turns toward its source. This reversion is not an
external journey. It is the orientation through which what proceeds
from the One becomes determinate as Intellect. The relation between
dependence and reversion is therefore constitutive of the intelligible
principle. It should not be reduced to a sequence of physical
movements or temporal stages.

\subsection{Soul: Mediation, Life, and the Sensible Order}
\label{subsec:ch08-plotinus-soul}

\phantomsection\label{para:VS.B01.P01.C08.S06.SS04.P001}
Soul is the third principal hypostasis. It proceeds from Intellect
and occupies a mediating position between intelligible reality and
the sensible cosmos. This mediation does not mean that Soul is merely
a boundary or a passive channel. Soul possesses its own activity and
is responsible for the articulation of life, motion, and order in
the sensible domain.

\phantomsection\label{para:VS.B01.P01.C08.S06.SS04.P002}
Plotinus distinguishes the primary activity of Soul, directed toward
Intellect, from its secondary activity, through which it gives order
to what lies below it. The distinction is not between two unrelated
souls or two temporally separated substances. It concerns two
directions of activity within the account of Soul. Its orientation
toward Intellect and its productive relation to the sensible order
belong to one metaphysical structure.

\phantomsection\label{para:VS.B01.P01.C08.S06.SS04.P003}
The sensible cosmos is not, for Plotinus, an independent reality
unrelated to intelligible principles. Its order is dependent upon
Soul, which communicates intelligible structure through formative
principles (\textit{logoi}). The sensible world is thus not simply
an accidental aggregate. It exhibits order, coordination, and
proportion because its constitution is related to intelligible
principles.

\phantomsection\label{para:VS.B01.P01.C08.S06.SS04.P004}
The activity of Soul is not adequately described as the physical
insertion of an immaterial object into a body. Plotinus repeatedly
resists spatial interpretations of the presence of Soul. Soul is
present to the body without being divided into spatial fragments
corresponding to bodily parts. Its unity is not distributed in the
same manner as the material components that it organizes.

\phantomsection\label{para:VS.B01.P01.C08.S06.SS04.P005}
The relation between Soul and the sensible world also requires a
distinction between the cosmic Soul and individual souls. The
cosmic order is not simply the sum of isolated organisms or private
psychic activities. Individual souls belong to a broader order of
life, while their particular activities remain distinguishable.
The unity of the cosmos therefore does not eliminate the plurality
of souls or the differences among living beings.

\subsection{Procession, Abiding, and Reversion}
\label{subsec:ch08-plotinus-three-moments}

\phantomsection\label{para:VS.B01.P01.C08.S06.SS05.P001}
A central pattern in Plotinus's metaphysics is commonly expressed
through three terms: abiding (\textit{monē}), procession
(\textit{proodos}), and reversion (\textit{epistrophē}). These terms
describe the dependence and activity of the hypostases. They should
not be read as three chronological episodes or as a mechanical
sequence of production.

\phantomsection\label{para:VS.B01.P01.C08.S06.SS05.P002}
Abiding indicates that a principle does not cease to be what it is
through the activity that proceeds from it. Procession indicates
the dependence of a derivative principle upon its source. Reversion
indicates the orientation of what proceeds toward its source, through
which it receives or realizes its determinate order. The three
moments are conceptually distinguishable but belong to one account
of ontological dependence.

\phantomsection\label{para:VS.B01.P01.C08.S06.SS05.P003}
The pattern is not uniform in every respect. The One, Intellect,
and Soul do not perform identical operations, nor does each possess
the same degree or kind of unity. The schema therefore provides a
general structure of dependence without erasing the differences
between the principles. A reading that treats the three moments as
interchangeable functions would obscure the architecture it is
intended to explain.

\phantomsection\label{para:VS.B01.P01.C08.S06.SS05.P004}
This dynamic account also prevents the hierarchy from being
misrepresented as a static ladder. Each principle is related both
to what is prior to it and to what follows from it. Its position
within the order is expressed through its distinctive mode of
activity, dependence, and orientation. The system is hierarchical,
but the hierarchy is not merely a ranking of isolated entities.

% ============================================================
% Project: The Philosophy of Viability
% Book: Book I — Ontology of Viability
% Book ID: VS.B01
% File: VS.B01.P01.C08.S06.SS05.P004.DG01.plotinian_procession_and_reversion.tex
% Asset: Diagram — Plotinian Procession and Reversion
% Purpose: Show ontological dependence without temporalizing it
% ============================================================

\begin{figure}[htbp]
    \centering
    \begin{tikzpicture}[
        node distance=1.35cm and 2.0cm,
        principle/.style={
            draw=ViabilityDarkBlue,
            rounded corners=3pt,
            line width=0.9pt,
            fill=ViabilityPaleBlue,
            text=ViabilityBlack,
            align=center,
            minimum width=3.2cm,
            minimum height=1.0cm,
            font=\small
        },
        relation/.style={
            -{Latex[length=2.5mm]},
            line width=0.9pt,
            draw=ViabilityGold
        },
        return/.style={
            -{Latex[length=2.5mm]},
            line width=0.9pt,
            dashed,
            draw=ViabilityGreen
        },
        annotation/.style={
            font=\footnotesize,
            text=ViabilityGray,
            align=center
        }
    ]

    \node[principle] (one) {
        \textbf{The One / Good}\\
        Absolute simplicity
    };

    \node[principle, below=of one] (nous) {
        \textbf{Intellect}\\
        Being and intelligible multiplicity
    };

    \node[principle, below=of nous] (soul) {
        \textbf{Soul}\\
        Life and mediation
    };

    \node[principle, below=of soul] (cosmos) {
        \textbf{Sensible Cosmos}\\
        Ordered, differentiated manifestation
    };

    \draw[relation] (one) -- node[right, annotation] {
        \textit{Procession}\\
        ontological dependence
    } (nous);

    \draw[relation] (nous) -- node[right, annotation] {
        \textit{Procession}\\
        derivative activity
    } (soul);

    \draw[relation] (soul) -- node[right, annotation] {
        \textit{Productive activity}\\
        formative order
    } (cosmos);

    \draw[return] (nous.east) .. controls +(1.8,0.2) and +(1.8,-0.2) ..
        node[right, annotation, xshift=0.65cm] {
            \textit{Reversion}\\
            toward the source
        } (one.east);

    \draw[return] (soul.east) .. controls +(2.2,0.2) and +(2.2,-0.2) ..
        node[right, annotation, xshift=1.05cm] {
            Contemplative\\
            orientation
        } (nous.east);

    \node[annotation, below=0.45cm of cosmos] {
        Schematic representation only: the relations are ontological,
        not temporal stages.
    };

    \end{tikzpicture}
    \caption{The Plotinian structure of procession and reversion.
    The diagram represents asymmetrical ontological dependence and
    return, not a temporal chain of events.}
    \label{fig:plotinian-procession-reversion}
\end{figure}

% END OF ASSET

\subsection{Matter, Formative Principles, and Sensible Manifestation}
\label{subsec:ch08-plotinus-matter-form}

\phantomsection\label{para:VS.B01.P01.C08.S06.SS06.P001}
Plotinus's account of sensible bodies retains some Aristotelian
vocabulary while altering its metaphysical function. Matter
(\textit{hulē}) and form remain important, but sensible matter is
not simply the potential constituent that becomes an actual
substance through the reception of form in the Aristotelian sense.
Plotinus's analysis assigns matter a more radically indeterminate
and privative status.

\phantomsection\label{para:VS.B01.P01.C08.S06.SS06.P002}
Sensible forms are related to intelligible principles through
formative principles or \textit{logoi}. These principles account
for the determinate characteristics through which sensible things
appear as the kinds of things they are. The sensible order is
therefore not self-explanatory. Its determinations are dependent
upon intelligible structures mediated through Soul.

\phantomsection\label{para:VS.B01.P01.C08.S06.SS06.P003}
The analogy between Plotinus's matter and a receptive substrate
must be handled carefully. Plotinian matter is not simply an
ordinary material component awaiting craftsmanship. It is described
as radically lacking determinate qualities and as incapable of
providing its own formal organization. The language of receptivity
does not mean that matter independently contributes a positive
principle of determination equal to the formative activity.

\phantomsection\label{para:VS.B01.P01.C08.S06.SS06.P004}
The relation between intelligible form and sensible appearance
is also not a relation of spatial transmission. The higher principle
does not travel downward as a body moves through space. The
sensible manifestation is dependent upon an intelligible order
without requiring that this order become spatially extended or
divided. The language of image, trace, and participation expresses
this asymmetrical relation of dependence.

\phantomsection\label{para:VS.B01.P01.C08.S06.SS06.P005}
This account yields a complex conception of sensible reality.
The sensible world is derivative, but it is not therefore simply
meaningless or wholly evil. Plotinus defends the beauty and
ordered character of the cosmos against views that treat embodiment
as an absolute fall from reality. The limitations of sensible
existence are connected with materiality and its privative character,
but the cosmos as an ordered whole remains an expression of
intelligible principles.

% ============================================================
% Project: The Philosophy of Viability
% Book: Book I — Ontology of Viability
% Book ID: VS.B01
% File: VS.B01.P01.C08.S06.SS06.P005.TL01.plotinian_ontological_principles.tex
% Asset: Table — Plotinian Ontological Principles
% Purpose: Compare the three hypostases and sensible reality
% ============================================================

\begin{table}[htbp]
    \centering
    \small
    \renewcommand{\arraystretch}{1.35}
    \begin{tabularx}{\textwidth}{
        >{\raggedright\arraybackslash}p{2.6cm}
        >{\raggedright\arraybackslash}X
        >{\raggedright\arraybackslash}X
    }
        \toprule
        \textbf{Principle} &
        \textbf{Ontological character} &
        \textbf{Explanatory role} \\
        \midrule

        The One / Good &
        Absolute simplicity; beyond determinate being and
        ordinary intellection. &
        First principle and source of the being and goodness
        of everything else. \\

        \addlinespace

        Intellect (\textit{Nous}) &
        Unified intelligible multiplicity; the Forms are
        internal to its living intelligibility. &
        The domain of Being, Forms, and determinate
        intelligibility. \\

        \addlinespace

        Soul (\textit{Psychē}) &
        A distinct hypostasis with an orientation toward
        Intellect and a productive relation to what follows. &
        Mediation of intelligible order into life, motion,
        and the sensible cosmos. \\

        \addlinespace

        Sensible Cosmos &
        Differentiated and extended reality, dependent upon
        formative principles and Soul. &
        The ordered sensible manifestation of a higher
        intelligible structure. \\

        \bottomrule
    \end{tabularx}
    \caption{The principal levels in Plotinus's metaphysical
    architecture. The table is schematic and does not imply
    that the principles are interchangeable or merely degrees
    of one homogeneous entity.}
    \label{tab:plotinian-ontological-principles}
\end{table}

% END OF ASSET

\subsection{Unity, Multiplicity, and the Status of Difference}
\label{subsec:ch08-plotinus-unity-difference}

\phantomsection\label{para:VS.B01.P01.C08.S06.SS07.P001}
Plotinus offers a particularly important account of the relation
between unity and multiplicity. The One is absolutely simple;
Intellect is a unified multiplicity; Soul mediates intelligible
order into a differentiated sensible cosmos; and bodies are
divisible and extended. These are not merely different quantities
of the same kind of unity. They are different modes of ontological
organization.

\phantomsection\label{para:VS.B01.P01.C08.S06.SS07.P002}
The unity of Intellect does not abolish the distinctness of Forms.
Its multiplicity is internal to an intelligible whole. Likewise,
the unity of the cosmos does not abolish the differences among
living beings, bodies, and activities. Plotinus therefore provides
a resource for distinguishing unity from numerical singularity
and from the aggregation of externally related parts.

\phantomsection\label{para:VS.B01.P01.C08.S06.SS07.P003}
Nevertheless, Plotinus's hierarchy assigns different degrees of
unity to different levels of reality. The One is the supreme
principle, and the lower levels are increasingly differentiated.
This is a distinctive metaphysical commitment. It should not be
transferred without argument into a framework in which the
ontological status of identity, plurality, and manifestation
remains independently under examination.

\phantomsection\label{para:VS.B01.P01.C08.S06.SS07.P004}
The distinction between the One and the intelligible multiplicity
also raises a question about the status of difference. If determinate
difference belongs to the intelligible order, then the first
principle cannot be described as one determinate identity among
others. Yet it does not follow that every account of identity must
be derived from a single absolutely simple source. That further
conclusion belongs to Plotinus's metaphysical architecture and
requires independent justification before it can be adopted elsewhere.

\subsection{Critical Examination of Plotinus's Explanatory Architecture}
\label{subsec:ch08-plotinus-critical-examination}

\phantomsection\label{para:VS.B01.P01.C08.S06.SS08.P001}
The explanatory strength of Plotinus's system lies in its attempt
to account for differentiated reality without treating the first
principle as one more differentiated item. The distinction between
the One, Intellect, and Soul allows him to articulate different
forms of unity, intelligibility, life, and sensible organization.
The system also avoids reducing sensible order to an accidental
collection of material events.

\phantomsection\label{para:VS.B01.P01.C08.S06.SS08.P002}
A first difficulty concerns the explanatory transition from absolute
simplicity to determinate multiplicity. Plotinus describes procession
and reversion as necessary features of the relation between principles,
but the precise force of this necessity remains philosophically
demanding. If the One is absolutely self-sufficient and beyond
determination, the account must explain why its causal productivity
does not compromise its simplicity, and why the resulting multiplicity
is not merely asserted as a consequence of the first principle.

\phantomsection\label{para:VS.B01.P01.C08.S06.SS08.P003}
A second difficulty concerns the asymmetry between the higher and
lower principles. The system explains lower reality through higher
principles, but the lower does not reciprocally constitute the higher.
This asymmetry is central to Plotinus's metaphysics. Its explanatory
success depends on whether the account can establish the relevant
dependence without presupposing the hierarchy it seeks to justify.

\phantomsection\label{para:VS.B01.P01.C08.S06.SS08.P004}
A third difficulty concerns the status of matter. The radical
privative characterization of matter provides a way to explain
deficiency and the limitations of sensible existence. Yet it raises
the question of whether a principle characterized primarily by
privation can adequately explain the determinate conditions under
which sensible manifestations occur. The distinction between
matter's privative status and the positive formative activity of
Soul must therefore remain explicit.

\phantomsection\label{para:VS.B01.P01.C08.S06.SS08.P005}
A fourth difficulty concerns the relation between the unity of the
cosmos and the plurality of its constituents. Plotinus gives a
substantial account of cosmic order and coordination. However,
the unity of an ordered whole does not by itself settle whether
the identities of its constituents are ontologically derivative,
independent, or dependent in more than one respect. The system's
answer is shaped by its own hierarchy and cannot be generalized
without further argument.

\subsection{Comparative Assessment from the Ontology of Viability}
\label{subsec:ch08-plotinus-comparative-assessment}

\phantomsection\label{para:VS.B01.P01.C08.S06.SS09.P001}
The comparison with the Ontology of Viability begins with a
methodological restriction: historical concepts must first be
understood within their own explanatory setting. Plotinus's One
is not to be renamed as the Transcendent Ground; Intellect is not
to be identified with Structural Determination; Soul is not to be
identified with a Hub; and sensible reality is not to be equated
with Material Manifestation. Such correspondences would require
arguments that the present comparison has not established.

\phantomsection\label{para:VS.B01.P01.C08.S06.SS09.P002}
The first point of comparison is the distinction between unity
and singularity. Plotinus's account of Intellect as a unified
multiplicity demonstrates that unity need not mean the absence
of internal distinction. This is relevant to the ontology of
identity because plurality of determinations, relations, or
manifestations should not be confused with plurality of identities.
The comparison, however, does not establish that Plotinus's
intelligible Forms and the identities examined in this book have
the same ontological status.

\phantomsection\label{para:VS.B01.P01.C08.S06.SS09.P003}
The second point concerns ontological dependence. Plotinus offers
a strongly asymmetric account in which the lower depends upon
the higher. The Ontology of Viability must examine dependence
without assuming that every relation of dependence forms a single
vertical hierarchy. A relation may be constitutive, enabling,
manifestational, or otherwise conditioned; its type and scope
must be specified before any priority claim is accepted.

\phantomsection\label{para:VS.B01.P01.C08.S06.SS09.P004}
The third point concerns manifestation. Plotinus's account of
images, traces, participation, and formative principles provides
a rich historical example of how sensible appearance may be
understood as dependent upon a non-sensible order. Yet the
Plotinian relation is not identical with the distinction between
identity and manifestation developed in this book. In particular,
the dependence of sensible appearance upon intelligible principles
does not establish that manifestation is merely an inferior copy
or that the identity of what manifests is exhausted by its
intelligible source.

\phantomsection\label{para:VS.B01.P01.C08.S06.SS09.P005}
The fourth point concerns the limits of conceptual predication.
Plotinus's treatment of the One warns against assuming that every
ontological principle must be describable through the same
determinations that characterize ordinary beings. This warning
is relevant to the examination of open ontological questions.
It does not, however, settle the nature or status of the
Transcendent Ground, the Throne-Domain, or Light within the
Ontology of Viability.

\phantomsection\label{para:VS.B01.P01.C08.S06.SS09.P006}
The fifth point concerns the relation between a system's
explanatory architecture and its foundational commitments.
Plotinus's account is coherent only within a set of commitments
about simplicity, causal dependence, intelligibility, and the
Good. A comparative analysis must therefore distinguish what
follows from those commitments from what can be independently
established. The presence of a powerful explanatory system is
not, by itself, proof that its first principles are necessary.

\subsection{Comparative Verdict and Unresolved Questions}
\label{subsec:ch08-plotinus-verdict}

\phantomsection\label{para:VS.B01.P01.C08.S06.SS10.P001}
Plotinus's enduring contribution is his systematic articulation
of the relation between absolute unity, intelligible multiplicity,
psychic life, and sensible manifestation. His account demonstrates
that a metaphysics of unity need not reduce every form of plurality
to an external aggregate, and that sensible reality can be treated
as ontologically dependent without being dismissed as wholly
unreal or devoid of order.

\phantomsection\label{para:VS.B01.P01.C08.S06.SS10.P002}
His principal limitation, from the standpoint of the present
comparative inquiry, is not that he fails to distinguish levels
of reality. It is that his explanatory architecture is governed
by a specific hierarchy whose first principle is absolutely
simple and whose derivative principles are ordered by degrees
of unity and dependence. The validity and reach of that hierarchy
cannot be presumed outside Plotinus's own system.

\phantomsection\label{para:VS.B01.P01.C08.S06.SS10.P003}
The comparison therefore yields three provisional conclusions.
First, unity must not be equated with numerical singularity or
with the absence of internal difference. Second, ontological
dependence must be specified rather than inferred from a general
hierarchical picture. Third, the relation between intelligible
determination and sensible manifestation requires a careful
account that does not collapse the manifested entity into the
principle through which it is manifested.

\phantomsection\label{para:VS.B01.P01.C08.S06.SS10.P004}
The following questions remain open: whether an absolutely simple
first principle is necessary for an adequate ontology; whether
the plurality of identities can be explained without reducing
them to a single source; whether manifestation must be understood
through a model of image and participation; and whether the
relation between unity and plurality admits forms of ontological
organization that are not arranged in a single hierarchy.

\phantomsection\label{para:VS.B01.P01.C08.S06.SS10.P005}
These questions are not resolved by the historical comparison.
They are retained as questions for the independent development
of the Ontology of Viability. Plotinus is thus neither adopted
as a foundation nor dismissed as an obsolete metaphysician.
He remains a demanding interlocutor whose account clarifies
the explanatory power and the commitments of a strongly
hierarchical metaphysics.

\subsection*{Epistemic and Comparative Status of the Present Section}

The historical claims in this section are presented as
interpretive reconstructions supported by the cited primary
texts and specialist scholarship. The critical assessments are
arguments about the explanatory structure of Plotinus's system,
not claims that every interpretation of Plotinus is settled.
The comparative conclusions are provisional and do not establish
a direct identity between Plotinian principles and the categories,
hubs, or open ontological questions of the Philosophy of Viability.

\begin{itemize}
    \item \textbf{Historical reconstruction:} conditionally supported
    by the primary texts and specialist scholarship.
    \item \textbf{Interpretation of procession and reversion:}
    argued, with interpretive qualifications.
    \item \textbf{Critical assessment of the One-to-multiplicity
    transition:} argued as an open philosophical difficulty.
    \item \textbf{Comparative conclusions:} provisional and
    restricted to the distinctions explicitly established above.
    \item \textbf{Identification of Plotinian principles with
    Book I entities:} not established and not adopted.
\end{itemize}

\subsection*{Question for Reflection}

Does the unity of an intelligible whole require a single absolutely
simple source, or can unity and plurality be ontologically related
without being arranged in one exhaustive hierarchy? In answering,
distinguish numerical singularity, internal multiplicity, ontological
dependence, and the identity of what manifests.

% ============================================================
% END OF SECTION 06 OF CHAPTER 08
% ============================================================